\begin{explanationbox}
	\textbf{\textcolor{CExplanation}{一句话直觉}}
	指数乘积式与对数负值式可以互相转化，核心是指数与对数互为逆运算。

	\medskip
	\textbf{\textcolor{CExplanation}{核心拆解}}
	\begin{description}[style=nextline,leftmargin=2.8em,labelsep=0.5em,itemsep=0.45em]
		\item[\textbf{要点一（等价变形）：}]
			等式两边取自然对数，将乘积转化为加法。
		\item[\textbf{要点二（条件约束）：}]
			$x+a>0$ 保证对数的真数大于零，等价关系才能成立。
	\end{description}

	\medskip
	\textbf{\textcolor{CExplanation}{几何本质（函数交点）}}
	将等式改写为 $e^{x}=1/(x+a)$，则 $x$ 是 $y=e^{x}$ 与 $y=1/(x+a)$ 图像交点的横坐标。由于 $x+a>0$，有 $1/(x+a)>0$，故交点的纵坐标为正。

	\medskip
	\textbf{\textcolor{CExplanation}{代数意义（降维工具）}}
	通过取对数，将指数乘积式 $e^{x}(x+a)=1$ 化为线性式 $x+\ln(x+a)=0$，降低函数阶数，便于求导或消元。

	\medskip
	\textbf{\textcolor{CExplanation}{考点价值（隐零点换元）}}
	在导数综合题中，当极值点无法显式求解时，利用此关系可将指数或对数表达式替换为另一种形式，简化后续运算。
	\begin{itemize}[leftmargin=*,itemsep=0.5em,topsep=0.4em]
		\item \textbf{考法一（极值代换）：}证明中遇到 $e^{x}(x+a)=1$ 隐零点条件，代入 $\ln(x+a)=-x$ 消去对数。
		\item \textbf{考法二（不等式放缩）：}利用该关系实现指数与对数的相互转化，构造对称形式。
	\end{itemize}

	\medskip
	\textbf{\textcolor{CExplanation}{顿悟点}}
	当指数函数与一次式相乘等于常数时，立刻想到取对数，将指数“拉下来”变成加法，从而转化为 $x+\ln(x+a)=0$ 的形式。

	\medskip
	\textbf{\textcolor{CExplanation}{使用场景}}
	\begin{itemize}[leftmargin=*,itemsep=0.5em,topsep=0.4em]
		\item \textbf{场景一（隐零点代换）：}函数极值点处满足 $e^{x}(x+a)=1$，无法解出 $x$，代换为 $\ln(x+a)=-x$ 后简化运算。
		\item \textbf{场景二（不等式证明）：}需要比较 $e^{x}$ 与 $\ln x$ 时，通过构造该关系实现底数转换。
	\end{itemize}

\end{explanationbox}
